\section{An explicit radian-angle bridge using finite sums}
\label{sec:finite-trig}

The circle parameters $c,s$ in the algebraic theorem can be replaced by a
real angle $\theta$ with $\cos\theta=c$ and $\sin\theta=s$.  Here is a
construction with primitive recursive precision bounds.  It does not invoke
the fundamental theorem of calculus.  Absolute values of complex numbers
are shorthand for the elementary Euclidean norm; all the estimates below
can equivalently be checked after squaring and clearing positive rational
denominators.

\subsection{Sine and cosine from finite Taylor polynomials}

Put
\[
 T_m(t)=\sum_{k=0}^{m}\frac{(\mathrm i t)^k}{k!}.
\]
For an integer $B\geq1$, $|t|\leq B$, and $m\geq2B$, the tail bound is
\[
 \sum_{k=m+1}^{\infty}\frac{B^k}{k!}
 \leq 2\frac{B^{m+1}}{(m+1)!}
 \leq 2^{\,2B^2+2B-m}.                                      \tag{T1}
\]
Indeed, after index $2B$ successive terms have ratio at most $1/2$,
$B^{2B}\leq2^{2B^2}$, and a finite geometric-sum bound gives the displayed
inequality before taking any limit.  Thus
$m=2B^2+2B+n+4$ already makes the rational tail at most $2^{-n-4}$.
The resulting rapid Cauchy name defines
\[
 E(t)=\cos t+\mathrm i\sin t.
\]

For completeness, this definition also gives a uniform algorithm when
$t$ is supplied by a real name, not a rational number.  Suppose
$|t|\leq M$, where $M\geq1$ is an integer, and set $B=M+1$.
Evaluate the input name at precision
$n+2B^2+2B+8$, obtaining a rational $q$; evaluate $T_m(q)$ with
$m=2B^2+2B+n+6$.  The error in a polynomial value caused by changing
$t$ to $q$ is at most
\[
 |t-q|\sum_{j\geq0}\frac{B^j}{j!}
 \leq |t-q|\,2^{2B^2+2B+2}.
\]
This and (T1) give an error smaller than $2^{-n}$.  Factorials, finite
sums, and the displayed precision transformers are primitive recursive.

We record the elementary identities needed below.  The finite Cauchy
product and the binomial theorem give
\[
 E(t+u)=E(t)E(u),\qquad E(-t)=\overline{E(t)},\qquad E(0)=1.   \tag{T2}
\]
There is no appeal to a theorem about infinite series hidden here: if
$|t|,|u|\leq B$, all terms of total degree at most $m$ agree exactly in
$T_m(t)T_m(u)$ and $T_m(t+u)$.  The absolute sum of the remaining terms is
at most the tail in (T1) with $2B$ in place of $B$.  For example,
$m=8B^2+4B+n+4$ suffices: using the preceding bound on the exponential
sum, the errors from replacing each factor by its Taylor polynomial,
the finite Cauchy-product remainder, and the tail for $t+u$ together
are at most $3\cdot2^{-n-4}$.  This proves (T2) at any requested
precision.  Consequently
$E(t)\overline{E(t)}=1$, so $|E(t)|=1$.

For $|a|\leq1$, the same finite-series estimates give
\[
 |E(a)-1|\leq 2|a|,\qquad
 |E(a)-1-\mathrm ia|\leq |a|^2,                             \tag{T3}
\]
because $\sum_{k\geq1}1/k!\leq2$ and
$\sum_{k\geq2}1/k!\leq1$.  Factoring $E(h)$ as a product of $K$ copies
of $E(h/K)$, with a supplied integer $K\geq\max(1,|h|)$, and telescoping
the product yields the convenient uniform bound
\[
 |E(a)-E(b)|\leq 3|a-b|.                                   \tag{T4}
\]
The harmless constant $3$ leaves slack for rational approximations.

\subsection{Arctangent from dyadic rational sums}

For a positive integer $N$, define
\[
 A_N(t)=\frac{t}{N}\sum_{j=0}^{N-1}
                  \frac{1}{1+(tj/N)^2}.                    \tag{T5}
\]
These are finite rational expressions for rational $t$; all denominators
are at least one.  Suppose $|t|\leq M$, $M\geq1$.  The exact rational
identity for $f(x)=1/(1+x^2)$ gives, for $|x|,|y|\leq M$,
\[
 |f(x)-f(y)|
 =\frac{|x-y|\,|x+y|}{(1+x^2)(1+y^2)}
 \leq2M|x-y|.
\]
Write $h=t/N$.  Their difference is exactly
\[
 A_N(t)-A_{2N}(t)
 =\frac h2\sum_{j<N}\bigl(f(jh)-f(jh+h/2)\bigr).
\]
The preceding bound therefore gives
\[
 |A_N(t)-A_{2N}(t)|\leq\frac{M^3}{2N}.
\]
Summing the finite geometric progression proves that the dyadic sums
define a rapid Cauchy real $A(t)$, with
\[
 |A(t)-A_N(t)|\leq\frac{M^3}{N}
 \quad\hbox{when $N$ is a power of two}.                    \tag{T6}
\]
The notation $A$ is the usual arctangent: (T5) is also a direct
constructive definition by rational integration.  No integration theorem
is required for the identities used here.

Input error is especially easy to control.  For $0\leq r\leq1$, put
$g(t,r)=t/(1+t^2r^2)$.  Direct subtraction gives
\[
 g(t,r)-g(u,r)
 =\frac{(t-u)(1-tu r^2)}{(1+t^2r^2)(1+u^2r^2)},
\]
and $(t^2+u^2)r^2\geq2|tu|r^2$ shows that its absolute value is at most
$|t-u|$.  Hence every $A_N$, and then $A$, is $1$-Lipschitz.
For a real input satisfying $|t|\leq M$, obtain a rational $q$ at
precision $n+3$, set $B=M+1$, and use the dyadic mesh
\[
 N=2^{n+3B+4}.
\]
Since $B^3\leq2^{3B}$, (T6) and the input-error estimate give an
approximation to $A(t)$ with error less than $2^{-n}$.  This is an
explicit primitive recursive name transformer.

\subsection{The Cayley identity without differentiation}

Define the unit complex rational function
\[
 w(x)=\frac{1+\mathrm ix}{1-\mathrm ix}
     =\frac{1-x^2+2\mathrm ix}{1+x^2}.
\]
We prove
\[
                     E(2A(t))=w(t).                        \tag{T7}
\]
Fix a dyadic $N\geq2M$, put $h=t/N$, and $x_j=jh$.
Because $x_jh\geq0$, the number
$D_j=1+x_j^2+x_jh$ is at least one.  Exact multiplication gives
\[
 \frac{w(x_j+h)}{w(x_j)}
       =\frac{D_j+\mathrm ih}{D_j-\mathrm ih}.
\]
Writing $r_j=h/D_j$, the difference of the right side from
$1+2\mathrm ir_j$ has norm at most $2h^2$.
Also
\[
 \left|r_j-\frac{h}{1+x_j^2}\right|
 =\frac{|x_j|h^2}{D_j(1+x_j^2)}\leq Mh^2.
\]
Thus its difference from $1+2\mathrm i h/(1+x_j^2)$ is at most
$2(M+1)h^2$.  Since $|2h/(1+x_j^2)|\leq1$, (T3) shows that
\[
 \left|\frac{w(x_j+h)}{w(x_j)}
       -E\!\left(\frac{2h}{1+x_j^2}\right)\right|
 \leq(2M+6)h^2.                                            \tag{T8}
\]

All factors in the two products have norm one.  Telescoping the products
in (T8), using $w(0)=1$ and (T2), gives
\[
 |w(t)-E(2A_N(t))|
 \leq \frac{(2M+6)M^2}{N}.
\]
By (T4) and (T6),
\[
 |w(t)-E(2A(t))|
 \leq\frac{8M^3+6M^2}{N}
 \leq\frac{16(M+1)^3}{N}.                                  \tag{T9}
\]
For example, $N=2^{n+3(M+1)+5}$ makes this error at most $2^{-n-1}$
and ensures $N\geq2M$.  In the later uniform case $M=2$, the smaller
choice $N=2^{n+9}$ already makes it less than $2^{-n}$.
Thus an explicit finite estimate at any requested precision proves (T7).
This is a finite error argument throughout, not an application of the
fundamental theorem of calculus or an existence theorem for differential
equations.

Define $\pi=4A(1)$, the standard arctangent normalization of the radian
constant.  Equation (T7) gives $E(2A(1))=\mathrm i$, so (T2) gives
\[
                           E(\pi)=-1.                       \tag{T10}
\]

\subsection{A finite chart choice produces an angle}

Let $c^2+s^2=1$.  Compute once a rational $q$ with
$|q-c|\leq1/4$, and make the decidable rational comparison $q\geq0$.
If this holds, $1+c\geq3/4$; set
\[
             t=\frac{s}{1+c},\qquad \theta=2A(t).
\]
If $q<0$, then $1-c>3/4$; set
\[
             t=\frac{s}{1-c},\qquad \theta=\pi-2A(t).
\]
In either case $|t|\leq4/3$, so one may use the uniform bound $M=2$
in all computations.  In the first chart, $w(t)=c+\mathrm is$; in the
second, $-\overline{w(t)}=c+\mathrm is$.  These follow by multiplying
through by $(1+c)^2+s^2=2(1+c)$ or
$(1-c)^2+s^2=2(1-c)$, respectively.  Equations (T2), (T7), and (T10)
therefore give $E(\theta)=c+\mathrm is$ in both charts.

This is a valid existential angle constructor.  Its selected angle may
depend on the supplied name and need not be a globally continuous
single-valued principal argument on the circle.  Both chart formulas
produce a correct angle wherever that chart is selected.  The chart is
chosen once at a fixed input precision and retained at every later
output precision.

\paragraph{Arithmetic audit.}
Every sum, product, dyadic precision, and number of iterations appearing
above has a displayed primitive recursive bound.  The finite binomial,
geometric-sum, and telescoping identities use quantifier-free induction
after rational denominators have been cleared.  The output precision is
a free natural parameter in those induction arguments.  Universal
quantification over precision is introduced after each finite estimate
has been proved; induction on a Cauchy statement containing such a
quantifier is not used.
